\documentclass[a4paper,10pt,landscape]{article}

\usepackage[czech]{babel}
\usepackage[T1]{fontenc}
\usepackage[utf8]{inputenc}
\usepackage{amsmath,amssymb}
\usepackage{geometry}
\usepackage{enumitem}

\geometry{a4paper,landscape,left=8mm,right=8mm,top=8mm,bottom=8mm}
\pagestyle{empty}
\setlength{\parindent}{0pt}
\setlength{\parskip}{0pt}
\setlist[enumerate]{leftmargin=*,itemsep=2pt,topsep=2pt}

% Uprav název testu; jednotlivé varianty jsou vždy A/B, C/D atd.
\newcommand{\TestTitle}{Geometrická posloupnost}
\newcommand{\TestHeader}[1]{%
    {\bfseries TEST -- \TestTitle}\hfill{\bfseries Varianta #1}\\[2mm]
    Jméno a příjmení: \rule{4.5cm}{0.4pt}\hfill Třída: \rule{1.5cm}{0.4pt}\\[2mm]
    Datum: \rule{2.4cm}{0.4pt}\hfill Počet bodů: \rule{1.4cm}{0.4pt}\hfill Známka: \rule{1.4cm}{0.4pt}\\[1mm]
    \hrule\vspace{3mm}
}
\newcommand{\AnswerSpace}[1]{\vspace{#1}}

\begin{document}
\noindent
\begin{minipage}[t]{0.48\textwidth}
    \small
    \TestHeader{E}

    % ZADÁNÍ VARIANTY E
    \begin{enumerate}
        \item Rozhodněte, zda posloupnost $a_n=(-2)^{n+1}$ je geometrická. Pokud ano, určete $a_1$ a $q$.
        \AnswerSpace{19mm}
        \item Jsou dány členy $a_3=12$, $a_6=96$. Určete $q$, $a_1$ a $a_9$.
        \AnswerSpace{33mm}
        \item Určete $S_{10}$ pro geometrickou posloupnost s $a_1=6$, $q=-\frac12$.
        \AnswerSpace{34mm}
        \item Mezi čísla $2$ a $162$ vložte tři kladná čísla tak, aby všech pět čísel tvořilo geometrickou posloupnost.
        \AnswerSpace{29mm}
    \end{enumerate}
\end{minipage}%
\hfill\vrule width 0.4pt\hfill
\begin{minipage}[t]{0.48\textwidth}
    \small
    \TestHeader{F}

    % ZADÁNÍ VARIANTY F
    \begin{enumerate}
        \item Rozhodněte, zda posloupnost $a_n=\frac{5}{2^n}$ je geometrická. Pokud ano, určete $a_1$ a $q$.
        \AnswerSpace{19mm}
        \item Jsou dány členy $a_2=-6$, $a_5=48$. Určete $q$, $a_1$ a $a_8$.
        \AnswerSpace{33mm}
        \item Určete $S_9$ pro geometrickou posloupnost s $a_1=3$, $q=-2$.
        \AnswerSpace{34mm}
        \item Mezi čísla $3$ a $96$ vložte čtyři kladná čísla tak, aby všech šest čísel tvořilo geometrickou posloupnost.
        \AnswerSpace{29mm}
    \end{enumerate}
\end{minipage}
\end{document}
